\documentclass[journal]{IEEEtran}

\usepackage[utf8]{inputenc}
\usepackage[T1]{fontenc}
\usepackage{newtxtext,newtxmath}
\usepackage{amsmath,bm}
\usepackage{mathtools}
\usepackage{booktabs}
\usepackage{siunitx}
\usepackage{cite}
\usepackage{microtype}
\usepackage[breaklinks=true,hidelinks]{hyperref}

% -----------------------------------------------------------------------
\title{Ocean Wave Spectral Metrics: A Survey of Definitions,\\
Physical Interpretation, and Historical Development}

\author{%
    \IEEEauthorblockN{Mikhail S.\ Grushinskiy}
    \IEEEauthorblockA{Independent Researcher,
        Bareboat Necessities Project,
        \texttt{mgrouch@users.github.com}}
    \thanks{Manuscript submitted April 2026.}
}

\begin{document}
    \maketitle

% -----------------------------------------------------------------------
    \begin{abstract}
        The statistical characterisation of random ocean wave fields has
        evolved over eight decades from empirical wave-height tables compiled
        for wartime ship routing into a rigorous spectral framework used today
        in offshore engineering, coastal management, climate monitoring, and
        marine safety.
        This survey collects and annotates the principal spectral metrics that
        have emerged from that development: the canonical spectral moments and
        their central-moment derivatives; significant wave height, wave period
        estimators, and steepness; spectral bandwidth under the competing
        definitions of Cartwright--Longuet-Higgins, Goda, Kuik, and others;
        regularity and narrowness parameters; nonlinear crest statistics from
        Rayleigh and Tayfun distributions; groupiness and the Benjamin--Feir
        Index; wave energy and power; Ursell number, wave age, and breaking
        indicators; shallow-water dispersion corrections and the
        Battjes--Janssen breaking model; extreme-value return levels by
        Weibull and peak-over-threshold methods; and human comfort indices
        standardised under ISO~2631.
        A bias-correction framework for inverse-frequency moments under noisy
        instantaneous frequency estimation is also presented.
        The survey is intended as a consolidated mathematical reference for
        practitioners implementing real-time oceanographic analysis in embedded
        or shipborne instrumentation.
    \end{abstract}

    \begin{IEEEkeywords}
        ocean wave spectra, spectral moments, significant wave height,
        wave period, spectral bandwidth, wave groups, Benjamin--Feir Index,
        breaking probability, motion sickness, extreme value analysis,
        shallow water, historical survey
    \end{IEEEkeywords}

% -----------------------------------------------------------------------
    \section{Introduction}

    \IEEEPARstart{T}{he} scientific measurement and statistical
    description of ocean waves has a history stretching back to the early
    twentieth century, when naval architects and harbour engineers
    first sought quantitative criteria for ship design and breakwater
    sizing.
    Early efforts were necessarily empirical: observers on vessels or coastal
    stations recorded the heights and periods of individual waves by eye or
    with primitive tide gauges, and the resulting datasets were reduced to
    simple averages and exceedance tables.
    The mathematical foundations that now underpin modern oceanographic
    practice were laid only during and after the Second World War, in
    response to urgent operational requirements for surf-zone landing
    forecasts and open-ocean ship routing.

    The decisive conceptual advance was the recognition, due largely to
    Sverdrup and Munk \cite{Sverdrup1947}, that the sea surface is not a
    superposition of a small number of regular sinusoidal waves but a
    continuous stochastic process whose energy is spread across a broad
    spectrum of frequencies and directions.
    Their 1947 manual on wave forecasting introduced the notion of
    significant wave height---defined initially as the average height of the
    highest one-third of waves as judged by a trained observer---and
    provided empirical fetch-duration diagrams that translated wind conditions
    into expected sea states.
    The significant wave height proved remarkably robust as an engineering
    design parameter and has remained the central descriptor of sea severity
    ever since, even as its definition was later placed on a rigorous
    spectral footing.

    The next transformation came from the statistical mechanics of random
    processes.
    In a landmark sequence of papers, Longuet-Higgins showed that the
    sea surface could be modelled as a superposition of many independent
    sinusoidal components with random phases, yielding a Gaussian surface
    elevation distribution \cite{LonguetHiggins1952}.
    From this foundation, Cartwright and Longuet-Higgins \cite{Cartwright1956}
    derived the full statistical distribution of wave maxima and demonstrated
    that crest heights follow the Rayleigh distribution in the narrowband limit.
    Their spectral bandwidth parameter $\varepsilon$, which measures the
    deviation of the crest-height distribution from Rayleigh, became the
    first precisely defined measure of spectral shape.

    The introduction of the directional wave spectrum $S(\omega,\theta)$ and
    the operational measurement of its integral moments from pitch-and-roll
    buoys and wave-rider systems during the 1960s and 1970s transformed wave
    statistics from a theoretical framework into a practical engineering tool.
    The JONSWAP field experiment \cite{Hasselmann1973} established empirical
    forms for the spectral peak shape and growth rates under fetch-limited
    conditions, providing the canonical JONSWAP spectral model that remains
    standard in engineering design codes today.
    Goda's extensive numerical experiments on wave statistics with spectral
    simulation \cite{Goda1970, Goda1985} produced the comprehensive tables
    and empirical formulas for characteristic heights, periods, and
    exceedance probabilities that are still cited in harbour and coastal
    structure design.

    Parallel development of nonlinear corrections to the Gaussian--Rayleigh
    framework addressed the systematic underestimation of large-crest
    probabilities in steep seas.
    Longuet-Higgins \cite{LonguetHiggins1963} computed the effect of
    second-order bound harmonics on the statistical distributions, and
    Tayfun \cite{Tayfun1980} derived a compact and practically useful
    second-order correction to the Rayleigh crest exceedance probability
    that is now widely used in offshore structural reliability.

    The instability of deep-water wave trains to sideband perturbations,
    discovered theoretically by Benjamin and Feir \cite{BenjaminFeir1967}
    and later connected to the formation of extreme or rogue waves, added a
    further layer of complexity to wave statistics in steep and
    long-crested seas.
    Quantitative assessment of this instability through the Benjamin--Feir
    Index became an element of rogue-wave risk guidelines in the 1990s and
    2000s.

    More recently, shallow-water and coastal applications have motivated
    extensions of the spectral framework to account for depth-dependent
    dispersion, wave shoaling, refraction, and the probabilistic breaking
    model of Battjes and Janssen \cite{Battjes1978}, while the expansion of
    recreational and commercial offshore activities has elevated human
    comfort and motion-sickness prediction---standardised through
    ISO~2631 \cite{ISO2631}---to a routine design consideration.

    This survey collects the complete set of scalar metrics that have emerged
    from this history, provides their mathematical definitions referenced to
    primary sources, and comments on their physical meaning and the
    conditions under which each is most informative.
    The presentation is organised by metric family: spectral moments and
    central moments (Sections~\ref{sec:moments}--\ref{sec:central}),
    frequency and bandwidth parameters
    (Sections~\ref{sec:freq}--\ref{sec:regularity}), wave periods
    (Section~\ref{sec:periods}), heights and steepness
    (Section~\ref{sec:heights}), crest probabilities and extremes
    (Sections~\ref{sec:probabilities}--\ref{sec:extremes}), groupiness
    and instability (Section~\ref{sec:groupiness}), energy and power
    (Section~\ref{sec:energy}), nonlinear and development metrics
    (Section~\ref{sec:nonlinear}), breaking (Section~\ref{sec:breaking}),
    comfort (Section~\ref{sec:comfort}), shallow-water extensions
    (Section~\ref{sec:shallow}), environmental coupling
    (Section~\ref{sec:env}), data quality diagnostics
    (Section~\ref{sec:quality}), and bias correction under noisy
    frequency estimates (Section~\ref{sec:bias}).

% -----------------------------------------------------------------------
    \section{Spectral Moments}
    \label{sec:moments}

    The frequency-domain one-sided wave spectrum $S(\omega)$ encodes the
    distribution of wave energy across angular frequency $\omega > 0$.
    In the directional case, $S(\omega)$ is the integral of the directional
    spectrum $S(\omega,\theta)$ over all directions $\theta$.
    The $n$-th spectral moment is defined as
    \begin{equation}
        m_n = \int_0^\infty \omega^n\,S(\omega)\,\mathrm{d}\omega,
        \qquad n \in \mathbb{Z}.
        \label{eq:moments_integral}
    \end{equation}
    In a time-domain context where the instantaneous frequency $\omega(t)$
    and displacement envelope power $P_\eta(t)$ are tracked continuously,
    \eqref{eq:moments_integral} is approximated by the running expectation
    \begin{equation}
        M_n = \bigl\langle P_\eta(t)\;\omega(t)^n \bigr\rangle,
        \label{eq:moments_running}
    \end{equation}
    which coincides with $m_n$ whenever the narrowband assumption holds
    \cite{Ochi1998}.
    In practice, moments up to order four and down to order $-1$ are
    routinely computed; higher orders amplify high-frequency noise
    while lower orders depend heavily on the spectral tail at low
    frequencies where measurement accuracy diminishes.

    The physical roles of each order are:
    \begin{itemize}
        \item $m_0$: total variance of surface displacement; proportional to the
        mean potential energy per unit area of the wave field.
        The identity $H_s \approx 4\sqrt{m_0}$ follows from the Rayleigh
        amplitude distribution \cite{Cartwright1956}.
        \item $m_1$: energy-weighted frequency; its ratio with $m_0$ gives the
        mean angular frequency of the spectrum.
        \item $m_2$: energy-weighted mean-square frequency.
        By Rice's upcrossing formula \cite{Rice1944, LonguetHiggins1952},
        the expected zero-upcrossing rate is
        $\nu_\uparrow = (2\pi)^{-1}\sqrt{m_2/m_0}$.
        \item $m_3$, $m_4$: higher-order moments that control the spectral
        skewness, kurtosis, and peakedness diagnostics due to
        Ochi \cite{Ochi1998} and Goda \cite{Goda1970}.
        \item $m_{-1}$: the negative-order moment used in the energy flux
        period $T_{-1,0} = 2\pi m_{-1}/m_0$, a quantity that arises
        naturally in wave energy resource assessment because the wave
        power is proportional to the group velocity integrated over
        the spectrum.
    \end{itemize}

% -----------------------------------------------------------------------
    \section{Central Moments and Spectral Shape}
    \label{sec:central}

    The shape of a spectral distribution is most naturally described
    relative to its mean rather than the origin.
    Letting $\bar\omega = m_1/m_0$ denote the energy-weighted mean frequency,
    the second, third, and fourth central spectral moments are
    \begin{align}
        \mu_2 &= \frac{m_2}{m_0} - \bar\omega^2,
        \label{eq:mu2}\\
        \mu_3 &= \frac{m_3}{m_0} - 3\bar\omega\frac{m_2}{m_0} + 2\bar\omega^3,
        \label{eq:mu3}\\
        \mu_4 &= \frac{m_4}{m_0}
        - 4\bar\omega\frac{m_3}{m_0}
        + 6\bar\omega^2\frac{m_2}{m_0}
        - 3\bar\omega^4.
        \label{eq:mu4}
    \end{align}
    These are the spectral analogues of the statistical moments of a
    probability distribution, with $S(\omega)/m_0$ playing the role of
    the probability density.

    $\mu_2$ measures spectral spread: it is small for a sharp swell peak
    and large for a broad wind sea.
    $\mu_3$ quantifies asymmetry of the spectrum about $\bar\omega$:
    a positive value indicates a heavier high-frequency shoulder, typical
    of actively wind-forced conditions.
    $\mu_4$ measures peakedness relative to a Gaussian spectral shape.

    Three normalised shape diagnostics follow:
    \begin{equation}
        \text{Skewness} = \frac{\mu_3}{\mu_2^{3/2}},
        \quad
        \text{Kurtosis} = \frac{\mu_4}{\mu_2^2},
        \quad
        \text{Excess kurtosis} = \frac{\mu_4}{\mu_2^2} - 3.
    \end{equation}

    \noindent
    Three closely related peakedness parameters appear in the engineering
    literature, all measuring the departure of the spectrum from a broad
    flat shape:
    \begin{itemize}
        \item \textbf{Ochi peakedness} \cite{Ochi1998}: $Q_p = m_0 m_4/m_2^2$.
        \item \textbf{Benassai parameter}: identical to $Q_p$; appears
        independently in Italian coastal engineering practice.
        \item \textbf{Goda peakedness} \cite{Goda1970}: $Q_G = 2(Q_p - 1)$.
    \end{itemize}
    All three are dimensionless and approach unity for a monochromatic sea.

% -----------------------------------------------------------------------
    \section{Frequency Metrics}
    \label{sec:freq}

    The mean angular frequency and its equivalent in hertz are
    \begin{equation}
        \bar\omega = \frac{m_1}{m_0}, \qquad \bar{f} = \frac{\bar\omega}{2\pi}.
    \end{equation}
    The mean frequency is the first and most fundamental descriptor of the
    ``typical'' oscillation rate of the sea surface.
    It is dominated by the location of the spectral peak in narrow spectra
    but is shifted toward higher frequencies relative to the peak in broad
    spectra where high-frequency energy is significant.

    The relative bandwidth (RBW) is the canonical dimensionless measure of
    spectral width normalised to the mean frequency:
    \begin{equation}
        \mathrm{RBW} = \frac{\sqrt{\mu_2}}{\bar\omega}.
        \label{eq:rbw}
    \end{equation}
    $\mathrm{RBW} \ll 0.1$ characterises a remote swell; values of
    $0.3$--$0.6$ are typical of a mixed wind-sea and swell condition;
    $\mathrm{RBW} > 0.6$ indicates a very broad or bimodal spectrum.

    An independent bandwidth estimate may be obtained from the
    \emph{phase-increment variance}, computed as the normalised variance
    of successive instantaneous phase increments $\Delta\phi/\Delta t$
    divided by $\bar\omega^2$.
    This estimator is insensitive to amplitude fluctuations and provides
    a consistency check against the moment-based RBW,
    particularly useful when envelope demodulation quality is uncertain.

% -----------------------------------------------------------------------
    \section{Spectral Bandwidth: Multiple Definitions}
    \label{sec:bandwidth}

    The concept of spectral bandwidth has attracted multiple competing
    definitions in the literature because the shape of the crest-height
    distribution depends on different combinations of moments depending
    on the order of the approximation.
    Four definitions are in common use:
    \begin{align}
        \varepsilon_\mathrm{CLH} &= \sqrt{1 - \frac{m_1^2}{m_0 m_2}}
        \qquad\text{\cite{Cartwright1956}},
        \label{eq:eCLH}\\
        \varepsilon_\mathrm{Goda} &= \sqrt{\frac{m_0 m_2}{m_1^2} - 1}
        \qquad\text{\cite{Goda1970}},
        \label{eq:eGoda}\\
        \varepsilon_\mathrm{Kuik} &= \frac{\sqrt{m_0 m_2 - m_1^2}}{m_1}
        \qquad\text{\cite{Kuik1988}},
        \label{eq:eKuik}\\
        \varepsilon_\mathrm{LH} &= \sqrt{\frac{m_0 m_2}{m_1^2} - 1}
        \qquad\text{(Longuet-Higgins)}.
        \label{eq:eLH}
    \end{align}

    The Cartwright--Longuet-Higgins parameter $\varepsilon_\mathrm{CLH}$
    was introduced in \cite{Cartwright1956} as the quantity whose square
    measures the deviation of the joint distribution of envelope and
    phase from the narrowband limit.
    It satisfies $0 \le \varepsilon_\mathrm{CLH} \le 1$, vanishes for a
    monochromatic wave, and equals unity for pure white noise.
    It is the most commonly cited bandwidth parameter in theoretical
    treatments.

    $\varepsilon_\mathrm{Goda}$ and $\varepsilon_\mathrm{LH}$ are algebraically
    identical and were introduced independently in engineering-oriented
    work on wave statistics.
    They are unbounded above and grow faster than $\varepsilon_\mathrm{CLH}$
    for very broad spectra.
    $\varepsilon_\mathrm{Kuik}$ \cite{Kuik1988} was derived from the
    analysis of pitch-and-roll buoy data and is the form recommended
    for general directional wave analysis; it reduces to
    $\varepsilon_\mathrm{Goda}$ when $m_1 \gg \sqrt{m_0 m_2 - m_1^2}$.

    In practice, the differences between the four definitions are small
    for the swell-dominated spectra common in open-ocean observations
    and become significant only for the broad or bimodal spectra of
    mixed sea states.

% -----------------------------------------------------------------------
    \section{Regularity and Narrowness}
    \label{sec:regularity}

    Related to bandwidth but distinct in emphasis are three
    \emph{regularity} descriptors.

    The \emph{spectral regularity index}
    \begin{equation}
        R_\mathrm{spec} = \exp(-\beta\cdot\mathrm{RBW})
    \end{equation}
    provides a normalised score in $[0,1]$ that decays exponentially with
    spectral width.
    The tuning constant $\beta$ is chosen empirically for a given
    application; $\beta = 3$ maps $\mathrm{RBW} = 0.3$ to
    $R_\mathrm{spec} \approx 0.4$, a useful mid-scale reference.

    The \emph{phase regularity}
    \begin{equation}
        R_\mathrm{phase} =
        \left\lVert
        \left\langle \frac{z}{|z|} \right\rangle
        \right\rVert,
    \end{equation}
    where $z(t)$ is the complex analytic signal, is the circular mean
    resultant length of the instantaneous phase; it equals unity for a
    perfectly coherent sinusoid and vanishes for uniformly distributed
    phase.
    It is sensitive to phase coherence rather than amplitude spread and
    therefore responds differently from moment-based bandwidths in
    mixed sea conditions.

    The classical \emph{spectral narrowness parameter} \cite{Cartwright1956}
    \begin{equation}
        \nu = \sqrt{\frac{m_0 m_2}{m_1^2} - 1}
        \label{eq:narrowness}
    \end{equation}
    is algebraically equivalent to $\varepsilon_\mathrm{Goda}$ and
    provides a compact alternative characterisation of spectral compactness
    that appears frequently in theoretical derivations of wave statistics.

% -----------------------------------------------------------------------
    \section{Wave Period Estimators}
    \label{sec:periods}

    The concept of a single representative wave period is inherently an
    approximation for a broadband random sea, and different choices of
    moment combination optimise different aspects of the period--spectrum
    relationship.
    The principal period estimators in common use are:
    \begin{align}
        T_z &= 2\pi\sqrt{\frac{m_0}{m_2}}
        \qquad\text{(mean zero-crossing period)},
        \label{eq:Tz}\\
        T_1 &= \frac{2\pi m_0}{m_1}
        \qquad\text{(energy-weighted mean period)},
        \label{eq:T1}\\
        T_{m02} &= 2\pi\sqrt{\frac{m_0}{m_2}}
        \qquad\text{(variance-based mean period)},
        \label{eq:Tm02}\\
        T_e &= \frac{2\pi m_{-1}}{m_0}
        \qquad\text{(energy flux period)},
        \label{eq:Te}\\
        T_{m0,-1} &= \frac{m_0}{m_{-1}},
        \quad
        T_{m1,-1} = \frac{m_1}{m_{-1}}.
        \label{eq:Tm01}
    \end{align}
    Note that $T_z$ and $T_{m02}$ are algebraically identical; both
    appear in the literature under different notational conventions.

    The mean zero-crossing period $T_z$ is the most widely used period
    in structural engineering design codes because it is directly related
    to the upcrossing rate via Rice's formula \cite{Rice1944}:
    \begin{equation}
        \nu_\uparrow = \frac{1}{T_{m02}} = \frac{1}{2\pi}\sqrt{\frac{m_2}{m_0}},
        \label{eq:upcrossing}
    \end{equation}
    giving the expected number of complete wave cycles per second.
    $T_e$ is preferred for wave energy applications because the wave group
    velocity, and hence the energy flux, is weighted toward lower frequencies
    than the energy-weighted mean.
    $T_1$ lies between $T_z$ and $T_e$ and is commonly used in DNV and
    ISO offshore design standards.

    The group period $T_g$---the mean spacing between wave groups as
    distinct from individual waves---is related to the second-order
    modulation of the wave envelope and enters the groupiness factor
    discussed in Section~\ref{sec:groupiness}.

% -----------------------------------------------------------------------
    \section{Significant Wave Height and Steepness}
    \label{sec:heights}

    \subsection{Significant Wave Height}

    The concept of significant wave height was introduced operationally
    by Sverdrup and Munk \cite{Sverdrup1947} as the average height of
    the highest one-third of waves in a record, closely matching the
    height subjectively estimated by experienced mariners.
    Cartwright and Longuet-Higgins \cite{Cartwright1956} showed that, for
    a Rayleigh amplitude distribution, this quantity equals $4\sigma$ where
    $\sigma = \sqrt{m_0}$ is the standard deviation of the surface
    elevation.
    The spectral definition
    \begin{equation}
        H_s = 4\sqrt{m_0}
        \label{eq:Hs}
    \end{equation}
    has since become the standard, replacing the earlier wave-by-wave
    definition in all modern spectral analysis frameworks.
    It is often written $H_{m0}$ to emphasise its origin in the zeroth
    spectral moment and to distinguish it from visually estimated heights.

    For a narrowband Rayleigh sea, the root-mean-square displacement
    \begin{equation}
        \sigma = \sqrt{m_0}
    \end{equation}
    is related to $H_s$ by $H_s = 2\sqrt{2}\,\sigma \approx 4\sigma$,
    with the factor $2\sqrt{2}$ arising from the mean of the upper-third
    of the Rayleigh distribution.

    \subsection{Characteristic Heights}

    The expected average of the highest $1/N$ fraction of waves in a
    Rayleigh sea is \cite{Longuet1963}
    \begin{equation}
        H_{1/N} = \sqrt{2}\,\sigma\!\left(\sqrt{2\ln N}
        + \frac{\gamma_\mathrm{E}}{\sqrt{2\ln N}}\right),
    \end{equation}
    where $\gamma_\mathrm{E} \approx 0.5772$ is the Euler--Mascheroni
    constant.
    For engineering practice, the values $H_{1/10}$ and $H_{1/100}$
    (the averages of the highest tenth and hundredth of waves respectively)
    are tabulated in design codes as multiples of $H_s$.

    \subsection{Wave Steepness}

    The deep-water wavelength associated with the mean zero-crossing period
    follows from the linear dispersion relation $\omega^2 = gk$:
    \begin{equation}
        L_0 = \frac{g T_z^2}{2\pi}.
    \end{equation}
    Wave steepness, the primary nonlinearity indicator in deep water, is
    \begin{equation}
        \mathrm{St} = \frac{H_s}{L_0}.
    \end{equation}
    Limiting steepness in deep water is $\mathrm{St}_\mathrm{max} \approx 0.142$
    (Miche criterion); whitecapping is commonly observed for
    $\mathrm{St} > 0.05$.
    Steepness enters the Tayfun nonlinear correction, the
    Benjamin--Feir Index, and the deep-water breaking probability proxy
    discussed in subsequent sections.

% -----------------------------------------------------------------------
    \section{Crest Probability Distributions}
    \label{sec:probabilities}

    \subsection{Rayleigh Distribution}

    For a stationary, ergodic, narrowband Gaussian sea surface, the
    amplitude of the wave envelope---and hence individual crest
    heights---follows the Rayleigh distribution.
    Cartwright and Longuet-Higgins \cite{Cartwright1956} established this
    result rigorously; the survivor function is
    \begin{equation}
        P(H_c > h) = \exp\!\left(-\frac{h^2}{2\sigma^2}\right).
        \label{eq:rayleigh}
    \end{equation}
    The Rayleigh model is accurate for well-developed, unidirectional seas
    but becomes unconservative (underestimates large-crest probability) as
    steepness increases, directional spreading narrows, or the spectrum
    broadens significantly.

    \subsection{Tayfun Nonlinear Correction}

    For steep near-breaking seas, second-order bound harmonics systematically
    raise individual crest heights above the linear prediction.
    Tayfun \cite{Tayfun1980} derived a compact second-order correction:
    \begin{equation}
        P(H_c > h) \approx
        \exp\!\left(-\frac{h^2}{2\sigma^2}\right)
        \exp\!\left(\Lambda\frac{h^3}{\sigma^3}\right),
    \end{equation}
    where $\Lambda$ is proportional to the wave steepness $\mathrm{St}$.
    This model reproduces laboratory and field observations of the
    enhanced probability of extreme crests significantly better than the
    pure Rayleigh model and is recommended for offshore structural
    reliability assessments in steep-sea environments.

    \subsection{Expected Extreme Crest Height}

    The expected maximum crest in a record containing $N$ waves follows
    the Gumbel asymptote of the Rayleigh distribution:
    \begin{equation}
        \mathbb{E}[H_\mathrm{max}] \approx
        \sigma\!\left(\sqrt{2\ln N}
        + \frac{\gamma_\mathrm{E}}{\sqrt{2\ln N}}\right),
        \label{eq:maxcrest}
    \end{equation}
    where $N = \nu_\uparrow T$ for a sea state of duration $T$.
    This expression grows slowly (logarithmically) with record length and
    provides a conservative estimate of the design crest height for a
    given exposure.

% -----------------------------------------------------------------------
    \section{Extreme Value Analysis}
    \label{sec:extremes}

    \subsection{Weibull Return Heights}

    For long-return-period design, particularly for offshore platform air
    gap and riser fatigue calculations, block maxima of the crest height
    distribution are fitted with a two-parameter Weibull model.
    Given shape parameter $\alpha$ and scale parameter $\beta$, the
    return height for return period $T_\mathrm{ret}$ satisfies
    \begin{equation}
        H_{T_\mathrm{ret}} = \beta
        \!\left(\ln\frac{T_\mathrm{ret}}{T_\mathrm{obs}}\right)^{1/\alpha},
    \end{equation}
    where $T_\mathrm{obs}$ is the observation duration.
    The Weibull model accommodates the non-Rayleigh tail behaviour
    observed in steep or shallow-water spectra.

    \subsection{Peak-Over-Threshold Analysis}

    The mean excess function above a threshold $u$ provides a complementary
    extreme-value diagnostic \cite{Coles2001}:
    \begin{equation}
        e(u) = \mathbb{E}[H_c - u \mid H_c > u].
    \end{equation}
    A linearly increasing $e(u)$ is consistent with a generalised Pareto
    exceedance tail; the slope yields the tail shape parameter.
    This metric is especially informative for storm-peak identification,
    rogue-wave assessment, and mooring fatigue analysis where individual
    exceedance events rather than bulk statistics are the primary concern.

% -----------------------------------------------------------------------
    \section{Wave Groupiness and Modulational Instability}
    \label{sec:groupiness}

    \subsection{Historical Background}

    The tendency of ocean waves to travel in groups---sequences of several
    large waves separated by intervals of smaller waves---was noted by
    mariners long before its theoretical basis was understood.
    Longuet-Higgins and Stewart \cite{LonguetHigginsStewart1964} analysed
    the interaction of wave groups with the mean water level through the
    concept of radiation stress, establishing that wave groups generate
    long-wave setdown and setup that drives nearshore circulation.
    The groupiness of a sea state is directly related to the spectral
    bandwidth: a narrow spectrum produces pronounced, regularly spaced
    groups, while a broad spectrum yields irregular, rapidly modulated
    wave envelopes.

    \subsection{Groupiness Factor}

    The groupiness factor
    \begin{equation}
        G = \frac{T_g}{T_z}
    \end{equation}
    measures the ratio of the average group period $T_g$ to the mean
    zero-crossing period $T_z$.
    Large values of $G$ indicate pronounced, slow-moving groups and are
    associated with enhanced fatigue loading on offshore structures and
    with the resonant excitation of moored vessels.

    \subsection{Benjamin--Feir Index}

    The theoretical instability of a uniformly progressive Stokes wave to
    small sideband perturbations was demonstrated by Benjamin and Feir
    \cite{BenjaminFeir1967}.
    The instability growth rate depends on the ratio of the wave steepness
    to the spectral bandwidth; their product defines the
    Benjamin--Feir Index:
    \begin{equation}
        \mathrm{BFI} = \frac{\sqrt{2}\;H_s/L_0}{\Delta f / f_p},
    \end{equation}
    where $f_p$ is the spectral peak frequency and $\Delta f$ is the
    bandwidth.
    $\mathrm{BFI} > 1$ signals conditions susceptible to modulational
    instability and the generation of abnormally large waves through
    the focusing of energy from the sideband components onto the carrier.
    The BFI has been adopted in rogue-wave risk guidelines and extreme
    sea-state design codes.

% -----------------------------------------------------------------------
    \section{Wave Energy and Power}
    \label{sec:energy}

    The mean wave energy density per unit surface area in deep water is
    \begin{equation}
        E = \rho g\,m_0,
    \end{equation}
    where $\rho$ is seawater density, a result that follows directly
    from the linear theory of water waves and the equipartition of energy
    between potential and kinetic forms \cite{Whitham1974}.

    The energy flux period arises naturally from the spectral integral
    of the group velocity-weighted energy:
    \begin{equation}
        T_{e,\mathrm{flux}} = \frac{m_{-1}}{m_0}.
    \end{equation}

    Deep-water wave power per unit crest length follows from linear wave
    theory as
    \begin{equation}
        P = \frac{\rho g^2}{64\pi}\,H_s^2\,T_e.
        \label{eq:power}
    \end{equation}
    Equation~\eqref{eq:power} is the standard formula used in wave energy
    resource characterisation and converter performance specification
    \cite{EMEC2009}.
    It makes explicit the strong ($H_s^2$) dependence of available power
    on sea severity and the linear ($T_e$) dependence on period, which
    favours longer-period swell over short-period wind sea for energy
    harvesting.

% -----------------------------------------------------------------------
    \section{Nonlinear and Development Metrics}
    \label{sec:nonlinear}

    Three parameters characterise the degree of nonlinearity or development
    of the sea state beyond the linear Gaussian model.

    \subsection{Ursell Number}

    The Ursell number \cite{Ursell1953}
    \begin{equation}
        Ur = \frac{H L^2}{h^3}
    \end{equation}
    measures the competition between wave nonlinearity (amplitude $H$ relative
    to depth $h$) and frequency dispersion (wavelength $L$ relative to depth)
    in water of finite depth.
    $Ur \ll 1$ corresponds to dispersive linear waves; $Ur \gg 1$
    characterises weakly dispersive, strongly nonlinear shallow-water waves
    described by cnoidal wave theory.
    It is an essential parameter for selecting the appropriate wave model in
    coastal and nearshore analysis.

    \subsection{Spectral Nonlinearity Parameter}

    The dimensionless steepness at the spectral scale,
    \begin{equation}
        \mu = \tfrac{1}{2}H_s k_p,
    \end{equation}
    where $k_p = (2\pi/T_p)^2/g$ is the peak wavenumber, scales directly
    with the amplitude of Stokes-type bound harmonics and enters the
    Tayfun crest probability correction as $\Lambda \propto \mu$.
    It provides a compact single-number summary of the degree to which
    nonlinear wave--wave interactions distort the Gaussian sea-surface
    statistics.

    \subsection{Wave Age}

    The wave age $c_p/U_{10}$---the ratio of the spectral peak phase speed
    to the 10-metre wind speed---was introduced in the wave-forecasting
    context by Sverdrup and Munk \cite{Sverdrup1947} and formalised by
    Pierson and Moskowitz \cite{PiersonMoskowitz1964} as the parameter
    controlling the fetch-limited growth of the spectrum toward a
    fully-developed equilibrium.
    Values below unity indicate a young, actively forced wind sea in which
    the waves are slower than the wind and still extracting energy;
    values exceeding unity indicate mature swell that has outrun its
    generating wind field.

% -----------------------------------------------------------------------
    \section{Wave Breaking}
    \label{sec:breaking}

    \subsection{Depth-Limited Breaking}

    In finite water depth, waves break when their height approaches a
    fixed fraction of the local water depth.
    The classic depth-limited breaking criterion
    $H_b \approx 0.78\,h$ is due to McCowan \cite{McCowan1894} for
    solitary waves and is widely applied to regular and irregular waves
    in engineering practice.
    A probabilistic breaking indicator consistent with the Rayleigh
    height distribution is
    \begin{equation}
        P_b^\mathrm{depth} = \exp\!\left(-\frac{(0.78\,h)^2}{2\sigma^2}\right),
    \end{equation}
    giving the fraction of wave crests in the statistical ensemble that
    would exceed the depth-limited maximum.
    The parameterisation of wave energy dissipation by random breaking
    through a probability density model was developed by Battjes and
    Janssen \cite{Battjes1978}, whose model remains the most widely used
    formulation in spectral wave models for depth-induced breaking.

    \subsection{Steepness-Limited Breaking}

    In deep water, breaking occurs when the wave slope reaches a critical
    value related to the Miche limiting steepness.
    A linear proxy commonly used in wave statistics is
    \begin{equation}
        P_b^\mathrm{steep} = 10\,\frac{H_s}{L_0},
    \end{equation}
    which grows toward unity as the sea approaches its deep-water breaking
    limit.
    Both depth-limited and steepness-limited indicators may be monitored
    simultaneously to distinguish shallow-water whitecapping from
    deep-water breaking in mixed observation environments.

% -----------------------------------------------------------------------
    \section{Human Comfort and Motion Sickness}
    \label{sec:comfort}

    The effect of ship motions on human comfort and performance became an
    active engineering concern with the expansion of offshore oil and gas
    operations in the 1970s and the growth of high-speed passenger ferries
    in the 1980s.
    The physiological response to whole-body vertical vibration is
    strongly frequency-dependent: the human vestibular system is most
    susceptible to accelerations in the range \SIrange{0.1}{0.5}{\hertz},
    which overlaps the frequency band of ocean swell.

    \subsection{Motion Sickness Dose Value}

    The ISO~2631 standard \cite{ISO2631} defines the Motion Sickness
    Dose Value (MSDV) as the root-mean-square of the frequency-weighted
    vertical acceleration integrated over the exposure duration $T$:
    \begin{equation}
        \mathrm{MSDV} = \sqrt{m_2}\,\sqrt{T}.
        \label{eq:msdv}
    \end{equation}
    The factor $\sqrt{m_2}$, the root-mean-square vertical acceleration,
    captures the spectral weighting relevant to the vestibular sensitivity
    band.
    MSDV has units of $\mathrm{m\,s^{-1.5}}$ and scales with both sea
    severity and exposure duration.

    \subsection{Derived Comfort Metrics}

    From MSDV the following application-level quantities are standardly
    derived:
    \begin{itemize}
        \item \textbf{Seasickness incidence:} the percentage of an average
        population expected to experience vomiting after a given
        exposure, computed from the ISO-calibrated logistic model
        as a function of MSDV and prevailing frequency.
        \item \textbf{Comfort rating:} a heuristic 0--100 scale, with 100
        representing flat-calm conditions and 0 severe discomfort,
        derived from a monotonic mapping of MSDV.
        \item \textbf{Vertical motion intensity} $\sqrt{m_2}$: expressed as
        RMS vertical acceleration, used for crew workload and
        equipment operability assessment.
        \item \textbf{Time to onset of discomfort:} an empirical estimate
        proportional to $(\mathrm{accel} \cdot f^{0.7})^{-1}$,
        approximating the latency before the average occupant
        experiences noticeable vestibular disturbance.
    \end{itemize}

% -----------------------------------------------------------------------
    \section{Shallow-Water Extensions}
    \label{sec:shallow}

    When observations are made in finite water depth $h$, the deep-water
    dispersion relation $\omega^2 = gk$ is replaced by the exact linear
    relation
    \begin{equation}
        \omega^2 = gk\tanh(kh),
        \label{eq:dispersion}
    \end{equation}
    which requires iterative solution for the wavenumber $k$ given $\omega$
    and $h$.
    The resulting depth-dependent quantities modify all period, wavelength,
    and energy flux estimates.

    Three further shallow-water metrics extend the standard open-ocean
    framework:
    \begin{itemize}
        \item \textbf{Shoaling coefficient $K_s$:} accounts for the change
        in wave amplitude as the group velocity $c_g = \mathrm{d}\omega/\mathrm{d}k$
        varies with depth.
        By energy flux conservation along a ray,
        $H_s \propto c_g^{-1/2}$, so shoaling amplifies waves as
        depth decreases.
        $K_s$ is used to relate nearshore observed heights to
        offshore reference conditions.
        \item \textbf{Refraction coefficient:} derived from Snell's law
        $c = \omega/k = \mathrm{const} \times \sin\theta$,
        it describes the rotation of wave crests toward depth
        contours as the phase speed decreases in shallower water.
        \item \textbf{Battjes--Janssen breaking probability
        \cite{Battjes1978}:} parameterises the dissipation of wave
        energy by random depth-induced breaking as the fraction of
        waves whose heights exceed $\gamma h$ (breaker index
        $\gamma \approx 0.78$), with the breaking probability
        computed from the truncated Rayleigh height distribution.
        This model is the foundation of the SWAN \cite{Booij1999}
        and WAVEWATCH~III \cite{Tolman2009} nearshore spectral
        wave models.
    \end{itemize}

% -----------------------------------------------------------------------
    \section{Environmental Coupling}
    \label{sec:env}

    \subsection{Wave Age Classification and Sea--Swell Partition}

    When wind speed $U_{10}$ is measured or estimated, the wave age
    $c_p/U_{10}$ enables categorical classification of the sea state:
    young ($c_p/U_{10} < 0.8$), mature ($0.8$--$1.2$), or
    swell-dominated ($> 1.2$).
    This classification connects directly to the Pierson--Moskowitz
    spectrum \cite{PiersonMoskowitz1964}, which corresponds to
    $c_p/U_{10} \approx 1.2$ (fully developed sea).
    The fraction of total spectral energy attributable to swell rather
    than locally generated wind sea is estimated from the wave age
    and informs ship routing, wave energy forecasting, and wave
    climate studies.

    \subsection{Radiation Stresses}

    The depth-integrated wave momentum flux tensor, or radiation stress,
    was introduced by Longuet-Higgins and Stewart
    \cite{LonguetHigginsStewart1964} to explain wave-driven phenomena
    such as set-down under wave groups, set-up in the surf zone, and
    longshore currents.
    The diagonal components are
    \begin{equation}
        S_{xx} \propto H_s^2, \qquad S_{yy} \propto H_s^2,
    \end{equation}
    with the proportionality constants depending on the ratio of group
    velocity to phase speed.
    Scalar proxies proportional to $H_s^2$ are used in engineering
    estimates of wave-driven coastal setup and sediment transport.

    \subsection{Bottom Orbital Velocity}

    The near-bed peak orbital velocity from linear wave theory is
    \begin{equation}
        u_b = \frac{\pi H_s}{T_z\,\sinh(kh)},
    \end{equation}
    which governs sediment mobilisation, scour beneath pipelines and
    monopile foundations, and benthic habitat disturbance.
    The $\sinh(kh)$ denominator ensures rapid decay with depth in
    deep-water conditions, becoming order unity near the breaking depth
    where bed stress is maximal.

    \subsection{Wave Energy Converter Capture Width}

    For wave energy harvesting applications the capture width ratio
    \begin{equation}
        \mathrm{CWR} = \frac{P_\mathrm{absorbed}}{P \cdot D}
    \end{equation}
    is the fraction of the incident wave power per unit crest length $P$
    that is extracted by a device of characteristic dimension $D$.
    The theoretical maximum for a heaving point absorber in deep water
    is $\mathrm{CWR}_\mathrm{max} = 1/(2\pi) \approx 0.16$ \cite{Evans1976}.

% -----------------------------------------------------------------------
    \section{Data Quality Diagnostics}
    \label{sec:quality}

    Spectral moment estimates computed from finite, possibly contaminated
    time series require accompanying quality flags for responsible use.

    The \emph{temporal stability index}---the running variance of $H_s$
    evaluated over a sliding reference window---quantifies the
    non-stationarity of the sea state.
    Large values indicate a rapidly evolving sea (e.g.\ storm onset or
    abatement) that may invalidate the stationarity assumptions underlying
    moment-based spectral estimators.

    The \emph{band-limited signal-to-noise ratio}, defined as the ratio of
    in-band variance to total variance for a user-specified frequency band,
    detects spectral energy outside the wave band of interest, including
    sensor noise floors, mechanical resonances, and parasitic contributions
    from vessel vibration.

    The \emph{data gap fraction}---the proportion of samples rejected as
    non-finite---provides a simple indicator of sensor health.
    Sustained gap fractions above a few percent indicate communication
    faults, ADC saturation, or accelerometer clipping that should be
    investigated before trusting derived spectral quantities.

% -----------------------------------------------------------------------
    \section{Bias Correction Under Noisy Frequency Estimation}
    \label{sec:bias}

    \subsection{Jensen-Inequality Bias}

    When instantaneous frequency estimates $\omega(t)$ carry measurement
    noise with variance $\sigma_\omega^2$, any spectral moment $M_n$ formed
    via \eqref{eq:moments_running} with $n < 0$ is biased upward relative
    to the true integral $m_n$ by Jensen's inequality, since
    $\omega \mapsto \omega^{-k}$ is strictly convex for $k > 0$:
    \begin{equation}
        \mathbb{E}\!\left[\omega^{-k}\right]
        > \frac{1}{\mathbb{E}[\omega]^k}.
    \end{equation}
    The same inflation propagates into all quantities derived from
    $m_{-1}$ and into ratios where a high positive power of $\omega$
    appears in the denominator.
    Under jittery frequency tracking---typical at low frequencies or
    during sea-state transitions---the uncorrected $H_s$ and $T_z$ may
    be overestimated by several percent, introducing spurious trends
    that could be misinterpreted as genuine sea-state changes.

    \subsection{First-Order Taylor Correction}

    Expanding $\omega^n$ around $\bar\omega = \mathbb{E}[\omega]$ to second
    order gives, for negative $n$,
    \begin{equation}
        \mathbb{E}\!\left[\omega^n\right] \approx
        \bar\omega^n
        \!\left(1 + c_n\,\frac{\sigma_\omega^2}{\bar\omega^2}\right),
        \label{eq:bias_expansion}
    \end{equation}
    with the coefficient $c_n = n(n-1)/2$ for integer $n$, yielding
    \begin{equation}
        c_{-4}=10,\quad c_{-3}=6,\quad c_{-2}=3,\quad c_{-1}=1,
        \quad c_0=0.
    \end{equation}
    Inverting \eqref{eq:bias_expansion} gives bias-corrected spectral moments:
    \begin{align}
        m_0^c &= \frac{m_0}{1 + 10\,\sigma_\omega^2/\bar\omega^2},
        \label{eq:M0c}\\
        m_1^c &= \frac{m_1}{1 + 6\,\sigma_\omega^2/\bar\omega^2},
        \label{eq:M1c}\\
        m_2^c &= \frac{m_2}{1 + 3\,\sigma_\omega^2/\bar\omega^2},
        \label{eq:M2c}\\
        m_3^c &= \frac{m_3}{1 + \sigma_\omega^2/\bar\omega^2},
        \label{eq:M3c}\\
        m_4^c &= m_4
        \quad\text{(no first-order correction required)},
        \label{eq:M4c}\\
        m_{-1}^c &= \frac{m_{-1}}{1 + 6\,\sigma_\omega^2/\bar\omega^2}.
        \label{eq:Mm1c}
    \end{align}

    All higher-level quantities---periods, heights, bandwidths, groupiness,
    energy flux, and comfort indices---should be computed from the
    corrected set $\{m_n^c\}$ when frequency noise is non-negligible.
    The correction is conservative: it reduces rather than inflates the
    estimates, providing a closer lower bound on the true spectral integrals.

    \subsection{Practical Effect}

    Bias correction is most beneficial for the significant wave height $H_s$,
    zero-crossing period $T_z$, and relative bandwidth RBW.
    It is especially important at long periods (swell below \SI{0.1}{\hertz})
    where the absolute frequency uncertainty is a larger fraction of
    $\bar\omega$, and during rapid sea-state transitions when the
    frequency tracker may exhibit transient jitter.
    For typical open-ocean conditions and well-converged trackers, the
    correction is of order $1$--$3\%$ and may be neglected; in challenging
    low-frequency or mixed-sea conditions it can reach $5$--$10\%$ and
    should be applied for any quantitative application.

% -----------------------------------------------------------------------
    \section{Summary of Metrics}
    \label{sec:summary}

    Table~\ref{tab:metrics} provides a consolidated reference of all
    spectral metrics discussed in this survey, with their symbols,
    moment-based definitions, and principal oceanographic applications.

    \begin{table*}[t]
        \centering
        \caption{Summary of ocean wave spectral metrics, symbols, and primary applications.}
        \label{tab:metrics}
        \renewcommand{\arraystretch}{1.22}
        \small
        \begin{tabular}{@{}llll@{}}
            \toprule
            \textbf{Metric} & \textbf{Symbol}
            & \textbf{Definition / source}
            & \textbf{Principal application} \\
            \midrule
            \multicolumn{4}{@{}l}{\emph{Spectral moments (integral basis)}} \\
            Zeroth moment   & $m_0$          & $\int S\,\mathrm{d}\omega$            & Wave energy; $H_s$ base. \\
            First moment    & $m_1$          & $\int \omega S\,\mathrm{d}\omega$     & Mean frequency numerator. \\
            Second moment   & $m_2$          & $\int \omega^2 S\,\mathrm{d}\omega$   & Upcrossing rate; steepness. \\
            Third moment    & $m_3$          & $\int \omega^3 S\,\mathrm{d}\omega$   & Spectral skewness. \\
            Fourth moment   & $m_4$          & $\int \omega^4 S\,\mathrm{d}\omega$   & Peakedness; kurtosis. \\
            Negative moment & $m_{-1}$       & $\int \omega^{-1} S\,\mathrm{d}\omega$& Energy flux; $T_e$. \\
            \midrule
            \multicolumn{4}{@{}l}{\emph{Frequency, bandwidth, and regularity}} \\
            Mean frequency      & $\bar\omega,\bar f$ & $m_1/m_0$               & Central frequency. \\
            Relative bandwidth  & RBW                 & $\sqrt{\mu_2}/\bar\omega$& Spectral width (dimensionless). \\
            CLH bandwidth       & $\varepsilon_\mathrm{CLH}$ & \eqref{eq:eCLH} & Crest-stat. bandwidth \cite{Cartwright1956}. \\
            Goda bandwidth      & $\varepsilon_\mathrm{Goda}$& \eqref{eq:eGoda}& Engineering spectral width \cite{Goda1970}. \\
            Kuik bandwidth      & $\varepsilon_\mathrm{Kuik}$& \eqref{eq:eKuik}& Directional buoy analysis \cite{Kuik1988}. \\
            Narrowness          & $\nu$               & \eqref{eq:narrowness}   & Compact bandwidth measure. \\
            Spectral regularity & $R_\mathrm{spec}$   & $e^{-\beta\,\mathrm{RBW}}$& Sinusoidality score. \\
            Phase regularity    & $R_\mathrm{phase}$  & Circular mean resultant & Phase coherence score. \\
            \midrule
            \multicolumn{4}{@{}l}{\emph{Wave periods}} \\
            Zero-crossing period  & $T_z,T_{m02}$       & \eqref{eq:Tz}           & Structural design; wave counting. \\
            Energy flux period    & $T_e$               & \eqref{eq:Te}           & Wave energy resource. \\
            Mean period           & $T_1$               & \eqref{eq:T1}           & DNV / ISO design codes. \\
            Aux.\ period ratios   & $T_{m0,-1},T_{m1,-1}$& \eqref{eq:Tm01}        & Multi-moment cross-checks. \\
            Zero-upcrossing rate  & $\nu_\uparrow$      & \eqref{eq:upcrossing}   & Wave counting; fatigue. \\
            \midrule
            \multicolumn{4}{@{}l}{\emph{Wave heights and steepness}} \\
            Significant wave height & $H_s (= H_{m0})$ & $4\sqrt{m_0}$           & Universal sea-state descriptor. \\
            RMS displacement        & $\sigma$           & $\sqrt{m_0}$            & Gaussian amplitude reference. \\
            Wave steepness          & St                 & $H_s/L_0$               & Nonlinearity; breaking onset. \\
            High-percentile crests  & $H_{1/10},H_{1/100}$& Rayleigh quantiles    & Structural peak-load design. \\
            Expected maximum crest  & $H_\mathrm{max}$   & \eqref{eq:maxcrest}     & Air-gap / freeboard design. \\
            \midrule
            \multicolumn{4}{@{}l}{\emph{Spectral shape and peakedness}} \\
            Skewness      & ---               & $\mu_3/\mu_2^{3/2}$           & Wind-sea vs.\ swell asymmetry. \\
            Kurtosis      & ---               & $\mu_4/\mu_2^2$               & Peakedness vs.\ Gaussian. \\
            Ochi peakedness & $Q_p$           & $m_0 m_4/m_2^2$               & Spectral sharpness \cite{Ochi1998}. \\
            Goda peakedness & $Q_G$           & $2(Q_p-1)$                    & Engineering peakedness \cite{Goda1970}. \\
            \midrule
            \multicolumn{4}{@{}l}{\emph{Crest probabilities and extreme values}} \\
            Rayleigh exceedance   & $P(H_c>h)$  & \eqref{eq:rayleigh}         & Linear crest statistics \cite{Cartwright1956}. \\
            Tayfun exceedance     & $P(H_c>h)$  & Nonlinear correction         & Steep-sea crests \cite{Tayfun1980}. \\
            Weibull return height & $H_{T_\mathrm{ret}}$ & Weibull fit of maxima & Long-return-period design. \\
            POT mean excess       & $e(u)$      & GPD tail                     & Storm-peak extremes. \\
            BFI                   & BFI         & Steepness / bandwidth        & Rogue-wave risk \cite{BenjaminFeir1967}. \\
            Groupiness factor     & $G$         & $T_g/T_z$                    & Group loading; mooring resonance. \\
            \midrule
            \multicolumn{4}{@{}l}{\emph{Energy, power, and development}} \\
            Wave energy density   & $E$         & $\rho g m_0$                 & Energy content per unit area. \\
            Wave power            & $P$         & \eqref{eq:power}             & WEC resource assessment. \\
            Ursell number         & $Ur$        & $HL^2/h^3$                   & Nonlinearity vs.\ dispersion \cite{Ursell1953}. \\
            Nonlinearity param.   & $\mu$       & $\frac{1}{2}H_s k_p$         & Bound-harmonic amplitude. \\
            Wave age              & $c_p/U_{10}$& Phase speed / wind speed     & Sea development stage \cite{PiersonMoskowitz1964}. \\
            \midrule
            \multicolumn{4}{@{}l}{\emph{Wave breaking}} \\
            Depth-limited break.\ prob. & $P_b^\mathrm{depth}$  & Rayleigh truncation & Nearshore breaker fraction. \\
            Steepness break.\ proxy     & $P_b^\mathrm{steep}$  & $10\,H_s/L_0$       & Deep-water whitecapping. \\
            Depth-limited crest height  & $H_b$                 & $0.78\,h$           & Maximum nearshore height. \\
            Battjes--Janssen prob.       & ---                  & \cite{Battjes1978}  & Spectral model dissipation. \\
            \midrule
            \multicolumn{4}{@{}l}{\emph{Human comfort (ISO~2631)}} \\
            MSDV                     & ---          & \eqref{eq:msdv}           & Whole-body vibration dose. \\
            Seasickness incidence    & ---          & ISO logistic model        & Crew/passenger welfare. \\
            Comfort rating           & ---          & Heuristic 0--100 scale    & Operability assessment. \\
            Vertical motion intensity& $\sqrt{m_2}$ & RMS accel.                & Equipment operability. \\
            Time to onset            & ---          & Empirical formula         & Exposure planning. \\
            \midrule
            \multicolumn{4}{@{}l}{\emph{Shallow-water and environmental}} \\
            Shoaling coefficient     & $K_s$        & Group velocity ratio      & Coastal wave transformation. \\
            Refraction coefficient   & ---          & Snell's law               & Directional wave bending. \\
            Battjes--Janssen model   & ---          & \cite{Battjes1978}        & Nearshore spectral model. \\
            Wave age class           & $c_p/U_{10}$ & Young / mature / swell    & Sea-state classification. \\
            Sea--swell partition     & ---          & Wave age threshold        & Swell fraction of energy. \\
            Radiation stress         & $S_{xx},S_{yy}$& \cite{LonguetHigginsStewart1964}& Coastal setup / currents. \\
            Bottom orbital velocity  & $u_b$        & Linear theory             & Sediment transport; scour. \\
            WEC capture width ratio  & CWR          & $P_\mathrm{abs}/(P\,D)$   & Wave energy converter perf. \\
            \bottomrule
        \end{tabular}
    \end{table*}

% -----------------------------------------------------------------------
    \section{Conclusion}

    The body of spectral metrics reviewed in this survey represents the
    accumulated output of eight decades of oceanographic research,
    beginning with the wartime wave-forecasting manuals of Sverdrup and
    Munk and continuing through the theoretical foundations laid by
    Longuet-Higgins, Cartwright, Goda, Tayfun, Benjamin and Feir, and
    Battjes and Janssen, to the modern probabilistic frameworks for
    extreme-value analysis and human comfort standardised in ISO~2631.

    Despite their diverse origins, the metrics share a common mathematical
    skeleton: all are expressible as functions of the spectral moments
    $m_n$, which in turn are integrals of the one-sided spectrum against
    polynomial or power-law frequency weights.
    This structural unity means that, once an accurate estimate of six
    moments ($m_{-1}$ through $m_4$) is in hand, the entire catalogue of
    wave parameters can be synthesised algebraically without any
    further spectral processing.

    Practical computation of these moments from instrumented platforms
    introduces two systematic error sources: the finite duration of the
    observation window (stationarity bias) and noise in the instantaneous
    frequency estimate (Jensen-inequality bias).
    Both are quantifiable and correctable at first order, as outlined in
    Section~\ref{sec:bias}, and their magnitude should be assessed
    whenever the derived metrics are used in structural design, safety
    classification, or climate-trend analysis.

% -----------------------------------------------------------------------
    \begin{thebibliography}{99}

        \bibitem{Sverdrup1947}
        H.~U. Sverdrup and W.~H. Munk,
        ``Wind, sea, and swell: Theory of relations for forecasting,''
        \emph{U.S. Navy Hydrographic Office Pub.\ 601}, Washington, D.C., 1947.

        \bibitem{LonguetHiggins1952}
        M.~S. Longuet-Higgins,
        ``On the statistical distribution of the heights of sea waves,''
        \emph{J.\ Marine Res.}, vol.~11, no.~3, pp.~245--266, 1952.

        \bibitem{Rice1944}
        S.~O. Rice,
        ``Mathematical analysis of random noise,''
        \emph{Bell Syst.\ Tech.\ J.}, vol.~23, pp.~282--332, 1944;
        vol.~24, pp.~46--156, 1945.

        \bibitem{Cartwright1956}
        D.~E. Cartwright and M.~S. Longuet-Higgins,
        ``The statistical distribution of the maxima of a random function,''
        \emph{Proc.\ R.\ Soc.\ Lond.\ A}, vol.~237, no.~1209,
        pp.~212--232, 1956.

        \bibitem{LonguetHiggins1963}
        M.~S. Longuet-Higgins,
        ``The effect of non-linearities on statistical distributions
        in the theory of sea waves,''
        \emph{J.\ Fluid Mech.}, vol.~17, no.~3, pp.~459--480, 1963.

        \bibitem{Longuet1963}
        M.~S. Longuet-Higgins,
        ``On the joint distribution of wave periods and amplitudes in a
        random wave field,''
        \emph{Proc.\ R.\ Soc.\ Lond.\ A}, vol.~389, pp.~241--258, 1983.

        \bibitem{LonguetHigginsStewart1964}
        M.~S. Longuet-Higgins and R.~W. Stewart,
        ``Radiation stresses in water waves; a physical discussion,
        with applications,''
        \emph{Deep-Sea Res.}, vol.~11, no.~4, pp.~529--562, 1964.

        \bibitem{BenjaminFeir1967}
        T.~B. Benjamin and J.~E. Feir,
        ``The disintegration of wave trains on deep water.
        Part~1. Theory,''
        \emph{J.\ Fluid Mech.}, vol.~27, no.~3, pp.~417--430, 1967.

        \bibitem{Hasselmann1973}
        K.~Hasselmann et al.,
        ``Measurements of wind-wave growth and swell decay during the Joint
        North Sea Wave Project (JONSWAP),''
        \emph{Dtsch.\ Hydrogr.\ Z.\ Suppl.}, vol.~A8, no.~12, 1973.

        \bibitem{Goda1970}
        Y.~Goda,
        ``Numerical experiments on wave statistics with spectral simulation,''
        \emph{Rep.\ Port Harbour Res.\ Inst.}, vol.~9, no.~3,
        pp.~3--57, 1970.

        \bibitem{Tayfun1980}
        M.~A. Tayfun,
        ``Narrow-band nonlinear sea waves,''
        \emph{J.\ Geophys.\ Res.}, vol.~85, no.~C3,
        pp.~1548--1552, 1980.

        \bibitem{Battjes1978}
        J.~A. Battjes and J.~P.~F.~M. Janssen,
        ``Energy loss and set-up due to breaking of random waves,''
        in \emph{Proc.\ 16th Int.\ Conf.\ Coastal Engineering},
        Hamburg, Germany, 1978, pp.~569--587.

        \bibitem{PiersonMoskowitz1964}
        W.~J. Pierson and L.~Moskowitz,
        ``A proposed spectral form for fully developed wind seas based on
        the similarity theory of S.~A. Kitaigorodskii,''
        \emph{J.\ Geophys.\ Res.}, vol.~69, no.~24,
        pp.~5181--5190, 1964.

        \bibitem{Kuik1988}
        A.~J. Kuik, G.~P. van Vledder, and L.~H. Holthuijsen,
        ``A method for the routine analysis of pitch-and-roll buoy wave data,''
        \emph{J.\ Phys.\ Oceanogr.}, vol.~18, pp.~1020--1034, 1988.

        \bibitem{Ochi1998}
        M.~K. Ochi,
        \emph{Ocean Waves: The Stochastic Approach}.
        Cambridge Univ.\ Press, 1998.

        \bibitem{Goda1985}
        Y.~Goda,
        \emph{Random Seas and Design of Maritime Structures}.
        Univ.\ of Tokyo Press, 1985.

        \bibitem{Whitham1974}
        G.~B. Whitham,
        \emph{Linear and Nonlinear Waves}.
        Wiley-Interscience, New York, 1974.

        \bibitem{Ursell1953}
        F.~Ursell,
        ``The long-wave paradox in the theory of gravity waves,''
        \emph{Proc.\ Cambridge Philos.\ Soc.}, vol.~49, pp.~685--694, 1953.

        \bibitem{McCowan1894}
        J.~McCowan,
        ``On the highest wave of permanent type,''
        \emph{Philos.\ Mag.}, vol.~38, pp.~351--358, 1894.

        \bibitem{ISO2631}
        International Organization for Standardization,
        \emph{ISO~2631-1: Mechanical Vibration and Shock---Evaluation of Human
        Exposure to Whole-Body Vibration}.
        ISO, Geneva, 1997.

        \bibitem{Evans1976}
        D.~V. Evans,
        ``A theory for wave-power absorption by oscillating bodies,''
        \emph{J.\ Fluid Mech.}, vol.~77, no.~1, pp.~1--25, 1976.

        \bibitem{Coles2001}
        S.~Coles,
        \emph{An Introduction to Statistical Modeling of Extreme Values}.
        Springer, London, 2001.

        \bibitem{Booij1999}
        N.~Booij, R.~C. Ris, and L.~H. Holthuijsen,
        ``A third-generation wave model for coastal regions. Part~I: Model
        description and validation,''
        \emph{J.\ Geophys.\ Res.}, vol.~104, no.~C4,
        pp.~7649--7666, 1999.

        \bibitem{Tolman2009}
        H.~L. Tolman,
        ``User manual and system documentation of WAVEWATCH~III version~3.14,''
        \emph{NOAA/NCEP Tech.\ Note}, 2009.

        \bibitem{EMEC2009}
        European Marine Energy Centre,
        \emph{Assessment of Wave Energy Resource}.
        EMEC, Orkney, 2009.

    \end{thebibliography}

\end{document}
